% Varied path diagram for the principle of least action / calculus of variations.
% Shows the true trajectory q(t) and a varied path q(t)+\delta q(t) with the
% same fixed endpoints q(t_1), q(t_2), plus the pointwise variation \delta q(t).
%
% Compile standalone with:  pdflatex 变分路径_tikz.tex
\documentclass[tikz,border=6pt]{standalone}
\usepackage{amsmath}
\usetikzlibrary{arrows.meta, calc}

\begin{document}
\begin{tikzpicture}[>=Stealth, thick, every node/.style={font=\fontsize{20}{24}\selectfont}
]

  % ---- geometry -----------------------------------------------------
  \def\tone{1.0}    % t1
  \def\ttwo{9.0}     % t2
  \def\qone{1.3}    % q(t1)
  \def\qtwo{7.8}     % q(t2)

  % ---- axes -----------------------------------------------------------
  \draw[blue!70!black, ->] (-0.6,0) -- (10.3,0) node[right] {$t$};
  \draw[blue!70!black, ->] (0,-0.6) -- (0,9.2) node[above, at end] {$q$};

  % ---- endpoint guide lines -------------------------------------------
  \draw[gray, dashed] (\tone,0) -- (\tone,\qone);
  \draw[gray, dashed] (0,\qone) -- (\tone,\qone);
  \draw[gray, dashed] (\ttwo,0) -- (\ttwo,\qtwo);
  \draw[gray, dashed] (0,\qtwo) -- (\ttwo,\qtwo);

  \node[gray, below] at (\tone,0) {$t_1$};
  \node[gray, below] at (\ttwo,0) {$t_2$};
  \node[gray, left]  at (0,\qone) {$q(t_1)$};
  \node[gray, left]  at (0,\qtwo) {$q(t_2)$};

  % ---- the two paths ----------------------------------------------------
  % Cubic Bezier segments keep the figure compatible with the editor's
  % lightweight renderer.  Both paths share the endpoints and midpoint.
  \draw[blue]
    (\tone,\qone)
      .. controls (2.4,1.3) and (3.6,2.0) .. (5,4.55)
      .. controls (6.4,7.1) and (7.6,7.8) .. (\ttwo,\qtwo);
  \draw[red, densely dashed]
    (\tone,\qone)
      .. controls (3.4,1.3) and (4.6,1.4) .. (5,4.55)
      .. controls (5.4,7.7) and (6.6,7.8) .. (\ttwo,\qtwo);

  \fill (\tone,\qone) circle (1.6pt);
  \fill (\ttwo,\qtwo) circle (1.6pt);

  % ---- \delta q(t) arrows between the two curves -------------------------
  \draw[green!45!black, densely dashed, -{Stealth[length=4pt]}]
    (3.85,2.6) -- (3.85,2);
  \draw[green!45!black, densely dashed, -{Stealth[length=4pt]}]
    (6.4,6.65) -- (6.4,7.29);

  \node[green!35!black] at (3,3.4) {$\delta q(t)$};

\end{tikzpicture}
\end{document}
